\documentclass[11pt]{article}
\usepackage[T1]{fontenc}
\usepackage{lmodern,amsmath,amssymb,amsthm,mathtools,booktabs}
\usepackage[margin=1.05in]{geometry}
\usepackage[hidelinks]{hyperref}
\newtheorem{theorem}{Theorem}[section]
\newtheorem{lemma}[theorem]{Lemma}
\newtheorem{proposition}[theorem]{Proposition}
\newtheorem{corollary}[theorem]{Corollary}
\newcommand{\R}{\mathbb R}
\newcommand{\Lk}{\mathcal L_{v,k}}
\newcommand{\code}[1]{{\small\nolinkurl{#1}}}
\title{Exact truncated moment boundaries\protect\\
for a Gaussian--cubic polynomial functional}
\author{Project Stardust}
\date{Reviewed editorial draft, October 29, 2026}
\begin{document}
\maketitle
\begin{abstract}
We determine positivity of polynomial squares through degree four and
realization by positive measures through degree nine for the
formal moment sequence with exponential generating series
$\exp(vz^2/2+kz^3/6)$. For $v>0$, put $r=k^2/v^3$.
Quadratic, cubic and quartic square tests give successively
\[
 r\leq2,\qquad 12-24r-9r^2\geq0,\qquad
 8-40r-45r^3\geq0.
\]
The last inequality is exactly the
criterion for matching moments through degree eight by a
positive probability measure. Matching the prescribed ninth
moment requires the strict inequality. At the quartic null
boundary a positive measure matches all eight moments, but
the null polynomial forces a different ninth moment. A
parallel obstruction occurs between degrees six and seven.
Finite Jacobi matrices construct the representing measures
and steer an arbitrary ninth moment in the strict region.
An attained sixth-moment bound gives a concrete low-degree
diagnostic. The methods are classical truncated moment
theory; the explicit parameter formulas are matched to a
frozen Lean development, with no priority claim.
\end{abstract}

\section{The polynomial functional and the main boundaries}
Let $v,k\in\R$. Define $m_0=1$ and, for $n\geq0$,
\begin{equation}\label{eq:recursion}
 m_{n+1}=vn\,m_{n-1}+\frac{k}{2}n(n-1)m_{n-2},
\end{equation}
where terms with negative indices are zero. Define the
linear functional $\Lk$ on $\R[x]$ by $\Lk(x^n)=m_n$.
The generating series identity
\begin{equation}\label{eq:formal}
 \sum_{n\geq0}m_n\frac{z^n}{n!}
       =\exp(vz^2/2+kz^3/6)
\end{equation}
is an identity of formal power series: differentiation
gives \eqref{eq:recursion} and the constant coefficient
is one. It is not the assumption that a positive law has
this moment-generating function.

Say that $\Lk$ passes the degree-$d$ square test when
$\Lk(p^2)\geq0$ for every real polynomial of degree at
most $d$. A Borel probability measure represents the
sequence through degree $N$ when $\int |x|^j\,d\mu<\infty$
and $\int x^j\,d\mu=m_j$ for $0\leq j\leq N$.
Introduce
\begin{equation}\label{eq:determinants}
 \begin{split}
 \delta&=2v^3-k^2,\\
 D&=12v^6-24v^3k^2-9k^4,\\
 Q&=288v^{10}-1440v^7k^2-1620vk^6.
 \end{split}
\end{equation}

\begin{theorem}\label{thm:classification}
For $v>0$ the degree-two, degree-three and degree-four
square tests hold exactly when $\delta\geq0$, $D\geq0$
and $Q\geq0$, respectively. A positive probability measure
represents the moments through degree eight exactly when
$Q\geq0$, and through degree nine exactly when $Q>0$.
All positive constructions may be chosen with finite support.
If $Q>0$, any real ninth moment can be prescribed while
keeping the first eight moments fixed.
\end{theorem}

The case $v=0$ is degenerate: a positive measure with the
first three prescribed moments must be $\delta_0$ and
$k=0$. If $v<0$, its second moment already excludes a
positive measure. Thus the full degree-eight criterion is
$(v,k)=(0,0)$ or $v>0,Q\geq0$; at degree nine replace
$Q\geq0$ by $Q>0$.

The issue at $Q=0$ is not failure of the preceding square
test. It is compatibility of the next moment with the
null polynomial. We give the positive constructions and
the incompatibility proof explicitly.

The general tools are classical. Curto and Fialkow's
odd truncated Hamburger criterion tests both positivity
and a range condition; their even criterion likewise
distinguishes singular data~\cite{curtofialkow}.
Here the contribution under examination is an explicit
evaluation for a two-parameter sequence and a verified
exposition of its boundaries. We do not claim a new
general moment theorem or established novelty for these formulas.

\section{Moment matrices and exact square tests}
The recursion gives
\begin{equation}\label{eq:table}
\begin{array}{c|l@{\qquad}c|l}
 j&m_j&j&m_j\\\hline
0&1&5&10vk\\
1&0&6&15v^3+10k^2\\
2&v&7&105v^2k\\
3&k&8&105v^4+280vk^2\\
4&3v^2&9&1260v^3k+280k^3
\end{array}
\end{equation}
Let $H_d=(m_{i+j})_{0\leq i,j\leq d}$. The square test
is exactly $H_d\succeq0$, since the coefficient vector
of $p$ has quadratic form $\Lk(p^2)$.

Use the monic centered basis
\begin{align*}
 U_0&=1,& U_1&=x,& U_2&=x^2-v,\\
 U_3&=x^3-3vx-k,&
 U_4&=x^4-6vx^2-4kx+3v^2.
\end{align*}
Each $U_j$, $j\geq1$, has zero functional mean. This
change of basis is triangular with diagonal one. On
$(U_1,U_2,U_3)$ the Gram matrix is
\begin{equation}\label{eq:G3}
 G_3=\begin{pmatrix}
 v&k&0\\ k&2v^2&6vk\\0&6vk&6v^3+9k^2
 \end{pmatrix},\qquad \det G_3=D.
\end{equation}
Adjoining $U_4$ gives
\begin{equation}\label{eq:G4}
 G_4=\begin{pmatrix}
 v&k&0&0\\
 k&2v^2&6vk&6k^2\\
 0&6vk&6v^3+9k^2&36v^2k\\
 0&6k^2&36v^2k&24v^4+144vk^2
 \end{pmatrix},\qquad \det G_4=Q.
\end{equation}
These identities follow by multiplying the displayed
polynomials and using \eqref{eq:table}. In particular the
leading determinants of $H_2,H_3,H_4$ are $\delta,D,Q$.

\begin{proposition}\label{prop:positivity}
For $v>0$ the three square-test criteria in
Theorem~\ref{thm:classification} hold.
\end{proposition}
\begin{proof}
For $p=c+bx+ax^2$, direct completion of squares yields
\[
 v\Lk(p^2)=v(c+va)^2+(vb+ka)^2+\delta a^2.
\]
This proves sufficiency of $\delta\geq0$.
If $\delta<0$, take $(c,b,a)=(-v^2,-k,v)$; its square
has value $v\delta<0$.

Necessity of $D\geq0$ and $Q\geq0$ follows because a
positive semidefinite matrix has nonnegative determinant.
For sufficiency put $r=k^2/v^3\geq0$. If $D\geq0$,
then $12-24r-9r^2\geq0$, which forces $r<1/2$ and
therefore $\delta>0$. The leading $2\times2$ block of
$G_3$ is positive definite. Its scalar Schur complement
is $D/\delta\geq0$, so $G_3\succeq0$.
If $Q\geq0$, its factorization
\begin{equation}\label{eq:normalize}
 Q=36v^{10}(8-40r-45r^3)
\end{equation}
implies $r\leq1/5$. Hence
$D/v^6\geq12-24/5-9/25=171/25>0$.
Thus $G_3$ is positive definite and the last Schur
complement in $G_4$ is $Q/D\geq0$. This proves
$G_4\succeq0$. The constant basis vector is orthogonal
to these centered blocks and has norm one, completing
the argument for $H_d$.
\end{proof}

This proof also identifies all positive pivots needed in
the representing matrices below. To display exact null
directions, define
\begin{align}
 R_3={}&6vk^2U_1-6v^2kU_2+\delta U_3,\label{eq:R3}\\
 R_4={}&-18k^3(10v^3-3k^2)U_1
       +18vk^2(10v^3-3k^2)U_2\notag\\
       &-72v^2k(v^3-k^2)U_3+D U_4.\label{eq:R4}
\end{align}
Multiplication by the Gram matrices shows that $R_3$
is orthogonal to every polynomial of degree at most two,
and $R_4$ to every polynomial of degree at most three.
Their square values are
\begin{equation}\label{eq:nullnorms}
 \Lk(R_3^2)=\delta D,\qquad \Lk(R_4^2)=DQ.
\end{equation}
When the preceding pivots are positive, the monic
orthogonal polynomials have squared norms
\begin{equation}\label{eq:norms}
 h_0=1,\qquad h_1=v,\qquad h_2=\delta/v,
 \qquad h_3=D/\delta,\qquad h_4=Q/D.
\end{equation}
For example, $P_2=U_2-(k/v)U_1$, $P_3=R_3/\delta$
and $P_4=R_4/D$. These computations can be checked from
the small Gram matrices without expanding large determinants.

\section{A finite spectral construction, including a null last pivot}
We spell out the finite construction to handle the
semidefinite boundary rather than invoke positivity alone.
The lemma is standard Jacobi-matrix moment theory.

\begin{lemma}\label{lem:jacobi}
Let $L$ be a real linear functional on polynomials of
degree at most $2d$, with $L(1)=1$. Suppose $H_{d-1}$
is positive definite and $H_d$ is positive semidefinite.
There is a real symmetric $(d+1)\times(d+1)$ matrix $J_f$,
depending on one real parameter $f$, for which
\[
 \langle e_0,J_f^n e_0\rangle=L(x^n),\qquad0\leq n\leq2d.
\]
Its spectral probability measure at $e_0$ represents
these moments. If $h_d$ is the last monic orthogonal
squared norm, then
\begin{equation}\label{eq:steergeneral}
 \langle e_0,J_f^{2d+1}e_0\rangle
 =\langle e_0,J_0^{2d+1}e_0\rangle+h_df.
\end{equation}
The construction remains valid when $h_d=0$.
\end{lemma}
\begin{proof}
Give $V_{d-1}=\R[x]_{\leq d-1}$ the positive definite
inner product $L(pq)$. Let $q_0=1,q_1,\ldots,q_{d-1}$
be its orthonormal polynomials with positive leading
coefficients. Orthogonal projection of $x^d$ onto
$V_{d-1}$ defines the monic residual $P_d$ and
$h_d=L(P_d^2)\geq0$.
For $p=p_{<d}+c x^d$ in $V_d$, define a coordinate map
into $\R^{d+1}$ by taking the coordinates of
$p-cP_d\in V_{d-1}$ and giving it last coordinate
$c\sqrt{h_d}$. Orthogonality proves
\begin{equation}\label{eq:isometry}
 \langle C(p),C(q)\rangle=L(pq),\qquad p,q\in V_d.
\end{equation}
When $h_d=0$ this is a map of the semidefinite quotient;
no division by $h_d$ is used.

Specify the first $d$ columns of $J_f$ by
$J_f C(p)=C(xp)$ for $p\in V_{d-1}$.
These columns are symmetric on their first $d$ entries,
because $L(xpq)=L(pxq)$. Fill the remaining off-diagonal
entries by symmetry, and set the last diagonal to $f$.
This defines a real symmetric matrix, including at $h_d=0$.
By induction $J_f^j e_0=C(x^j)$ for $0\leq j\leq d$:
only the first $d$ columns are used until the last step.
For $n\leq2d$ choose $i,j\leq d$ with $i+j=n$. Then
\[
 \langle e_0,J_f^n e_0\rangle
 =\langle J_f^i e_0,J_f^j e_0\rangle
 =L(x^{i+j})
\]
by \eqref{eq:isometry}.

Let $J_f=\sum_\lambda\lambda\Pi_\lambda$ be its finite
orthogonal spectral decomposition. The measure
\begin{equation}\label{eq:spectralmeasure}
 \mu_f=\sum_\lambda\|\Pi_\lambda e_0\|^2\delta_\lambda
\end{equation}
has nonnegative weights summing to one. Its $n$th moment
is exactly $\langle e_0,J_f^n e_0\rangle$, and finite
support supplies all absolute moments.

Finally $J_f-J_0=f e_d e_d^\mathsf T$, while
$J_f^d e_0=C(x^d)$ is independent of $f$ and has last
coordinate $\sqrt{h_d}$. Thus
\[
 \langle e_0,J_f^{2d+1}e_0\rangle
 =\langle C(x^d),J_f C(x^d)\rangle
 =\langle C(x^d),J_0 C(x^d)\rangle+h_df.
\]
This proves \eqref{eq:steergeneral}.
\end{proof}

The symmetric matrices in the lemma are tridiagonal.
Indeed if $i<j-1$, orthogonality gives
$L(xq_iq_j)=0$, because $xq_i$ has degree below $j$.
If $P_j$ is monic with norm $h_j>0$, the neighboring
entry is $\sqrt{h_j/h_{j-1}}$ by its leading coefficient.
At $h_d=0$, the last row and column decouple. The
spectral measure at $e_0$ then lives on the preceding
block; an arbitrary last diagonal introduces no mass.

\section{Eight moments and a freely chosen ninth moment}
Assume $v>0$ and $Q\geq0$. By
Proposition~\ref{prop:positivity}, $H_3$ is positive
definite and $H_4$ positive semidefinite.
Lemma~\ref{lem:jacobi} with $d=4$ gives a positive
probability measure representing all moments through eight.
Conversely such a measure makes every degree-four square
nonnegative, so $Q\geq0$. This proves the degree-eight
classification, including its boundary.

For direct reproduction, the matrix is
\begin{equation}\label{eq:explicitJ}
 J_f=\begin{pmatrix}
 0&b_1&0&0&0\\
 b_1&a_1&b_2&0&0\\
 0&b_2&a_2&b_3&0\\
 0&0&b_3&a_3&b_4\\
 0&0&0&b_4&f
 \end{pmatrix},
\end{equation}
where
\begin{align}
 b_1&=\sqrt v,& b_2&=\sqrt\delta/v,&
 b_3&=\sqrt{vD}/\delta,& b_4&=\sqrt{\delta Q}/D,\label{eq:offdiag}\\
 a_1&=k/v,&
 a_2&=\frac{k(4v^3+k^2)}{v\delta},&
 a_3&=\frac{72v^2k(v^3-k^2)}D-\frac{6v^2k}\delta.\label{eq:diag}
\end{align}
All denominators are strictly positive. The off-diagonal
entries follow from \eqref{eq:norms}; the diagonal entries
are $\Lk(xP_j^2)/h_j$ for $j=1,2,3$.
For example $\Lk(xP_1^2)=k$, and inserting
$P_2=x^2-(k/v)x-v$ gives
$\Lk(xP_2^2)=4vk+k^3/v^2$, which yields $a_2$.
The same substitution of $R_3/\delta$ gives $a_3$.
Thus \eqref{eq:explicitJ} is the concrete form of the
construction, not a matrix whose moment matching is assumed.

Write $s_9(f)=\int x^9\,d\mu_f$. The steering formula is
\begin{equation}\label{eq:ninthsteer}
 s_9(f)=s_9(0)+(Q/D)f.
\end{equation}
If $Q>0$, any target $s\in\R$ is obtained by
$f=(s-s_9(0))D/Q$. In particular one can choose
$s=m_9=1260v^3k+280k^3$. This proves the positive
degree-nine construction and the arbitrary-target assertion.
At $Q=0$, however, \eqref{eq:ninthsteer} has zero slope.
The next section shows that the prescribed $m_9$ is
incompatible with every positive representation, not
merely with this chosen family of matrices.

\section{The null boundary excludes the prescribed ninth moment}
We first record the general compatibility argument.

\begin{lemma}\label{lem:null}
Suppose a positive measure represents $L$ through degree
$2d+1$, and $p$ has degree at most $d$ with $L(p^2)=0$.
Then $L(xp^2)=0$.
\end{lemma}
\begin{proof}
Polynomial integrability follows from the assumed
absolute moments. Since $p^2\geq0$ and its integral is
zero, $p=0$ almost everywhere. Therefore $xp^2=0$
almost everywhere, and its integral is zero.
\end{proof}

For the particular polynomial $R_4$ of \eqref{eq:R4},
direct multiplication with the moments through degree
nine gives the exact identity
\begin{equation}\label{eq:cross4}
\begin{split}
 v\Lk(xR_4^2)={}&(-9k^5+48v^3k^3-108v^6k)Q\\
 &+864v^{10}k(52v^6-272v^3k^2+213k^4).
\end{split}
\end{equation}
This is a polynomial identity in $v,k$, verifiable from
the finite table; the accompanying symbolic checker
independently verifies its transcription.

Now assume $v>0,Q=0$. Its normalized equation is
$8-40r-45r^3=0$, where $r=k^2/v^3\geq0$.
It implies $r\leq1/5$, and $k\ne0$ since the constant
term is eight. On $0\leq r\leq1/5$ the polynomial
$52-272r+213r^2$ is decreasing, as its derivative
$-272+426r$ is negative. Its minimum on this interval is
$153/25>0$. Hence the second line of \eqref{eq:cross4}
is nonzero, and $\Lk(xR_4^2)\ne0$.
But \eqref{eq:nullnorms} gives $\Lk(R_4^2)=0$.
Lemma~\ref{lem:null} rules out a positive degree-nine
representation. Together with the preceding section,
this completes Theorem~\ref{thm:classification} for $v>0$.

The degenerate cases stated after the theorem also follow
without a moment-extension theorem. If $v=0$, a positive
representing measure has $x=0$ almost everywhere, so
$k=0$ and it must be $\delta_0$. Conversely
\eqref{eq:recursion} gives $m_j=0$ for every $j>0$
when $(v,k)=(0,0)$. Negative $v$ is impossible because
$\int x^2\,d\mu\geq0$.

The boundary is nonempty. The polynomial
$g(r)=8-40r-45r^3$ is strictly decreasing for $r\geq0$,
with $g(0)>0$ and $g(1/5)<0$. It has a unique positive
root $r_4$, and $g(1/6)=9/8>0$ shows
\begin{equation}\label{eq:r4}
 1/6<r_4<1/5.
\end{equation}
For $v=1$, $k=\sqrt{r_4}$ supplies an explicit algebraic
parameter at which all eight moments are realizable but
the prescribed ninth is not.

\section{The preceding six-to-seven obstruction}
The smaller null boundary illustrates the same mechanism
at a different truncation and is part of the same story.

\begin{proposition}\label{prop:sixseven}
If $v>0$ and $D=0$, a positive probability measure
represents the moments through degree six, but none
represents them through degree seven.
\end{proposition}
\begin{proof}
The equation $D=0$ gives $12-24r-9r^2=0$, with
\begin{equation}\label{eq:r3}
 r=r_3=(2\sqrt7-4)/3<1/2.
\end{equation}
Thus $\delta>0$, $H_2$ is positive definite and $H_3$
is positive semidefinite. Lemma~\ref{lem:jacobi} with
$d=3$ constructs the degree-six representation.
The last pivot $h_3=D/\delta$ is zero, so it suffices
to retain the preceding three-dimensional block.

The null cubic $R_3$ satisfies $\Lk(R_3^2)=\delta D=0$.
Its next cross moment is
\begin{equation}\label{eq:cross3}
 \Lk(xR_3^2)=-14v^2kD+24v^5k(10v^3-17k^2).
\end{equation}
Again this follows by polynomial multiplication and
\eqref{eq:table}. At $D=0$ one has $k\ne0$ and
$10v^3-17k^2=v^3(10-17r_3)>3v^3/2>0$.
So \eqref{eq:cross3} is nonzero, contradicting
Lemma~\ref{lem:null} for any degree-seven representation.
\end{proof}

There is also an open gap between the cubic and quartic
square tests. At $v=1,k=1/2$, one has $r=1/4$ and
\[
 D=87/16>0,\qquad
 Q/36=8-40/4-45/64=-173/64<0.
\]
Continuity preserves both strict signs on a neighborhood.
Thus even strict cubic square positivity does not
guarantee quartic square positivity or eight-moment
realizability. This example is a direct specialization
of the exact criteria, not an appeal to numerical eigenvalues.

\section{An attained sixth-moment bound}
The formal sequence can fail much earlier at particular
parameters. A sharp bound explains the obstruction at
$v=k=1$ and supplies an independent finite construction.

\begin{proposition}\label{prop:sharp}
If a probability measure has moments
\[
 m_0=1,\quad m_1=0,\quad m_2=1,\quad m_3=1,
 \quad m_4=3,\quad m_5=t,
\]
and finite sixth absolute moment, then
\begin{equation}\label{eq:sharp}
 m_6\geq t^2-8t+26.
\end{equation}
For every real $t$ equality is attained by a finitely
supported probability measure.
\end{proposition}
\begin{proof}
Put $p_t(x)=x^3+(4-t)x^2+(t-7)x+(t-5)$.
Expanding its square with the prescribed first five
moments gives
\[
 \int p_t(x)^2\,d\mu=m_6-(t^2-8t+26),
\]
which proves the lower bound.
For attainment, take the spectral measure at $e_0$ of
\[
 A_t=\begin{pmatrix}0&1&0\\1&1&1\\0&1&t-5\end{pmatrix}.
\]
Its first four coordinate-power vectors are
\[
 e_0,\quad A_t e_0=(0,1,0),\quad
 A_t^2e_0=(1,1,1),\quad A_t^3e_0=(1,3,t-4).
\]
Their inner products give moments $m_0,\ldots,m_4$
as stated, and
\[
 \langle A_t^2e_0,A_t^3e_0\rangle=t,\qquad
 \|A_t^3e_0\|^2=1+9+(t-4)^2=t^2-8t+26.
\]
The spectral measure is a probability measure with at
most three support points and all absolute moments.
Its square integral of $p_t$ is zero, so its support
is contained in the real zero set of $p_t$.
\end{proof}

For the prescribed sequence at $v=k=1$, $m_5=10$ and
$m_6=25$, whereas \eqref{eq:sharp} requires $m_6\geq46$.
Indeed $p_{10}=x^3-6x^2+3x+5$ satisfies
$\mathcal L_{1,1}(p_{10}^2)=-21$.
The degree-two square test still passes, since
$k^2=1\leq2v^3$. This contrasts a finite positive
truncation with the next genuinely incompatible data.

\section{Normalization, interpretation and exact formal scope}
The scale $r=k^2/v^3$ is intrinsic for $v>0$.
Replacing $x$ by $\sqrt v\,y$ transforms the formal
series to $\exp(z^2/2+\kappa z^3/6)$ with
$\kappa=k/v^{3/2}$. Changing $k$ to $-k$ is reflection
of the real variable. Thus the displayed boundaries depend
only on $r=\kappa^2$. The hierarchy is
\[
 0<r_4<1/5<r_3<1/2<2,
\]
where $r_4$ is the unique positive root of
$8-40r-45r^3$ and $r_3=(2\sqrt7-4)/3$.
The strict gaps show why matching a few moments cannot
justify interpreting the whole functional as expectation
under a positive law. Our assertions stop at the stated
finite truncations; no all-degree positivity classification
is needed or supplied by this paper.

The proof uses elementary Gram matrices, Schur complements,
finite spectral measures and null-polynomial compatibility.
These are standard truncated Hamburger tools, including
the distinction between positivity and extendibility at
singular data in Curto--Fialkow~\cite{curtofialkow}.
Their odd-case range condition is reflected concretely
in \eqref{eq:cross3} and \eqref{eq:cross4}; the strict
case admits the free last diagonal in
\eqref{eq:ninthsteer}. A targeted preliminary search did
not settle priority for the explicit two-parameter
boundaries. Novelty remains unresolved.

The frozen verified baseline is
\begin{center}\small\ttfamily
bab688bb6140cd6316c80680095a914ea0d346db.
\end{center}
The exact final declarations are in the namespaces
\begin{center}\small
\code{MixedCubicFourthDegree},\quad
\code{MixedCubicEighthRealization},\\
\code{MixedCubicNinthRealization},\quad
\code{MixedCubicQuarticBoundary},\\
\code{MixedCubicBoundaryQuartic},\quad
\code{MixedCubicSharpSixthMoment},
\end{center}
with prefix \code{OperatorDerivativeRewrite.Foundation}.
The theorem map records declaration lines and frozen source
hashes. The general finite spectral lemma is a human
reorganization of the explicitly verified Jacobi matrices
and steering identities; it is not represented as a newly
named formal theorem. The rational example and the
normalization discussion are direct deductions, identified
separately from the mapped source statements.

\begin{thebibliography}{9}
\bibitem{curtofialkow} R. E. Curto and L. A. Fialkow,
Recursiveness, positivity, and truncated moment problems,
\emph{Houston J. Math.} \textbf{17} (1991), no.~4,
603--635. In particular Theorems 3.1 and 3.9.
\url{https://www.math.uh.edu/~hjm/v017n4/0603CURTO.pdf}.
\end{thebibliography}
\end{document}
